
\section{Panoramica}
	Il progetto prevede la realizzazione di un applicativo che consenta la vendita di biglietti per eventi teatrali (come opere e concerti) e la divulgazione di informazioni a riguardo. 
	Il sistema si articolerà, per l’acquirente, in due sezioni: la prima conterrà la lista di eventi disponibili all'acquisto, la seconda conterrà maggiori informazioni riguardo le opere, incluse quelle non in programma al momento. 
	Per la biglietteria saranno disponibili sezioni per la gestione degli eventi e delle prenotazioni, mentre per l'amministratore saranno disponibili sezioni per la gestione delle informazioni riguardo le opere e della struttura posti del teatro. 
	Dall’intervista con i responsabili marketing e ufficio stampa di un ente pesarese operante nel settore sono emerse varie necessità. 
	L’interfaccia dovrà essere ottimizzata per favorire la semplicità del processo di acquisto del biglietto: un’eccessiva quantità di informazioni in primo accesso ridurrebbe l’intuitività e l’accessibilità. 
	Il prezzo di un biglietto è determinato dalla tipologia di posto scelto e dal tipo di evento, nonché da una lista di sconti applicabili in base a determinati criteri. 
	Una volta confermata la prenotazione, il totem digitale stamperà una ricevuta contenente l'identificativo che potrà essere usato per effettuare il pagamento in biglietteria. 
	Le voci delle informazioni dell’evento (direttore d’orchestra, interpreti, musicisti e direttori artistici e altri) dovranno essere totalmente flessibili viste le notevoli differenze fra i vari tipi di eventi. 
	Un caso particolare di eventi sono le opere, che possono essere messe in scena solo tramite una regia (una direzione artistica che accoppia alle musiche e al libretto dell’opera i costumi, fondali e altri elementi necessari alla rappresentazione). 
	Infine, nella sezione anagrafica, per ogni opera dovrà essere memorizzato e mostrato all’utente lo storico delle regie presenti e passate.
\newpage